\chapter{Informes}\label{cap:informes}

\paraque Genera un informe en PDF y CSV sobre el estado de los expedientes o
la actividad de cada litigante en un periodo.

\disponible Administrador (todo el despacho) y Litigante (sus expedientes
asignados).

\begin{table}[H]
  \centering
  \caption{Tipos de informe.}\label{tab:tipos-informe}
  \begin{tabularx}{\textwidth}{L{4.8cm} Y}
    \toprule
    \textbf{Informe} & \textbf{Qué incluye} \\
    \midrule
    \textbf{Estado y carga de expedientes}
      & Expedientes creados en el periodo, su estado y el litigante
        asignado. \\
    \textbf{Actividad por litigante}
      & Documentos cargados, actuaciones registradas y plazos atendidos por
        cada litigante en el periodo. \\
    \bottomrule
  \end{tabularx}
\end{table}

\pasos
\begin{enumerate}
  \item En el menú, haga clic en \textbf{Informes}.
  \item En \textbf{Solicitar informe}, elija el tipo de informe.
  \item Indique el periodo: \textbf{Creación desde (UTC)} y
        \textbf{Creación hasta (excluida, UTC)}, o bien
        \textbf{Actividad desde (UTC)} y
        \textbf{Actividad hasta (excluida, UTC)}, según el tipo. El primer
        día se incluye y el último no; el periodo máximo es de 366 días.
  \item Si lo desea, filtre por \textbf{Estado administrativo} y por
        litigante.
  \item Haga clic en \textbf{Generar informe}.
  \item En \textbf{Solicitudes y avisos}, espere a que el estado cambie a
        \textbf{Disponible} y haga clic en \textbf{Consultar informe}.
  \item Haga clic en \textbf{Descargar PDF} o \textbf{Descargar CSV}.
\end{enumerate}

\begin{figure}[H]
  \centering
  \capturapendiente{capturas/A-39-informes.png}%
    {Informes de expedientes con la solicitud y un informe Disponible con Descargar PDF y CSV · Litigante}
  \caption{Solicitud y descarga de un informe.}
  \label{fig:manual-informes}
\end{figure}

\resultado Descarga el informe en el formato elegido. El PDF y el CSV
provienen de la misma captura de datos.

\begin{itemize}
  \item El informe se genera en el servidor; no necesita mantener la
        pantalla abierta.
  \item Las fechas del periodo se interpretan en UTC.
  \item Si el informe es demasiado grande, aparece \enquote{El informe supera
        el límite permitido. Solicita otro con un periodo o filtro más
        reducido.}
\end{itemize}

\nota{El informe muestra el estado al momento de generarlo; no reconstruye
el estado pasado del expediente ni califica resultados jurídicos.}
